\documentclass[11pt]{article}

\usepackage[T1]{fontenc}
\usepackage[utf8]{inputenc}
\usepackage{lmodern}
\usepackage{microtype}
\usepackage[margin=1in]{geometry}
\usepackage{amsmath,amssymb,amsthm,mathtools}
\usepackage{booktabs,tabularx,array}
\usepackage{enumitem}
\usepackage{xcolor}
\usepackage{hyperref}
\usepackage[nameinlink,noabbrev]{cleveref}

\hypersetup{
  colorlinks=true,
  linkcolor=blue!45!black,
  citecolor=blue!45!black,
  urlcolor=blue!45!black
}

\newtheorem{theorem}{Theorem}[section]
\newtheorem{proposition}[theorem]{Proposition}
\newtheorem{corollary}[theorem]{Corollary}
\theoremstyle{definition}
\newtheorem{definition}[theorem]{Definition}
\newtheorem{example}[theorem]{Example}
\theoremstyle{remark}
\newtheorem{remark}[theorem]{Remark}

\newcommand{\SEL}{\operatorname{SEL}}
\newcommand{\dist}{\operatorname{dist}}
\newcommand{\Conf}{\operatorname{Conf}}
\newcommand{\id}{\operatorname{id}}
\newcommand{\N}{\mathbb{N}}
\newcolumntype{Y}{>{\raggedright\arraybackslash}X}

\title{\textbf{Selection Reachability in Projection-Based Dynamics:}\\
\large Conditional Undecidability, Robust Encoding, and Finite-Capacity
Limits in Modal Triplet Theory}
\author{Peter Nero}
\date{Version 2, July 2026}

\begin{document}
\maketitle

\begin{abstract}
Projection, stable records, and threshold events do not by themselves imply
algorithmic undecidability.  Undecidability enters only after a dynamical
realization is shown to simulate a universal machine uniformly for arbitrarily
long admissible runs.  This paper replaces an earlier generic claim by an
auditable reachability framework.  It imports, without duplicating, the
conditional two-counter-machine theorem of the revised MTT Program A2 and
proves four complementary results.  First, a finite projection system can have
noninjective descent, an absorbing record basin, and completely decidable
selection.  Second, undecidable reachability descends through a projection only
under an effective semiconjugacy, admissibility, input, and target-fidelity
contract.  Third, a finite-horizon robust event is decidable once a complete
positive-margin enclosure certificate is supplied.  Fourth, a compact carrier
cannot encode infinitely many distinct configurations with one uniform
positive decoding radius; unbounded robust computation therefore requires
shrinking margins, a non-totally-bounded carrier, a growing family of carriers,
or another explicitly verified mechanism.  We also prove that every finite
trajectory prefix may be computable while unbounded target reachability is
undecidable.  Thus reachability undecidability is not a theorem that every
trajectory is uncomputable, that no shortcut exists, or that projection creates
universal computation.  Current MTT has not yet supplied a selected physical
embedding satisfying the complete unbounded robustness contract, so its
physical selection-reachability claim remains conditional.
\end{abstract}

\section*{Revision note for Version 2}
\begin{description}[leftmargin=3.3cm,style=nextline]
\item[Supersedes.] Version 1, DOI
\href{https://doi.org/10.5281/zenodo.18255391}{10.5281/zenodo.18255391}.
\item[Reason.] Version 1 inferred universal computation and undecidable
selection from noninjective projection, record stability, basin structure, and
finite admissibility.  Its simulator construction assumed, rather than proved,
unbounded counters, exact zero tests, uniform instruction fidelity, and
arbitrarily long admissible persistence.
\item[Resolution.] Version 2 makes the universal-machine reduction conditional
on the complete Program A2 embedding contract, proves what a projection must
preserve before undecidability can descend, and separates finite-horizon
verification from unbounded reachability.
\item[Retained content.] The distinction between chaos and algorithmic
reachability, the use of a fixed universal counter machine, and the selection
event as a possible target predicate remain useful once their hypotheses are
typed.
\item[Remaining boundary.] No selected physical MTT carrier currently provides
the full input code, step map, zero test, unbounded admissible persistence,
target fidelity, and robust neighborhoods required for the conditional
undecidability theorem.
\end{description}

This note records the release delta and is not part of the abstract.

\tableofcontents

\section{The corrected question}

There are several mathematically different ways in which a physical model can
be difficult to predict.  A trajectory can be sensitive to initial data.  A
finite computation can be expensive.  A reduced state can omit memory needed
for autonomous evolution.  An unbounded reachability question can be
undecidable.  These statements do not imply one another.

Version 1 moved too quickly between them.  It treated a many-to-one projection
as if it created universal computation, interpreted finite record margins as
unlimited counter storage, and then described undecidable reachability as the
absence of an effective dynamical law.  Each step is false without additional
data.

The corrected question is narrower:
\begin{quote}
Under which explicit encoding and descent conditions does the halting problem
reduce to the occurrence of a selected event in an effective physical
dynamics?
\end{quote}
This is a reachability question.  It is not automatically a claim about all
trajectories, all protocols, all physical theories, or the computational cost
of finite prediction.

The revision has a deliberately layered role.  The MTT Program A2 paper owns
the conditional two-counter-machine reduction~\cite{NeroA2}.  The present
paper does not rename or re-prove that theorem.  It supplies the missing
projection-transfer analysis, capacity obstruction, finite-horizon certificate
boundary, and explanatory counterexamples.

\section{Four claims that must not be conflated}

We use the computable admissible presentation introduced in Program
A2~\cite{NeroA2}.  In abbreviated form it consists of
\[
  \mathfrak C=(S,\iota,F,A,H),
\]
where \(S\) is an effectively represented state set, \(\iota\) is a computable
input encoder, \(F:S\rightharpoonup S\) is a partial computable step map with an
effective domain, \(A\subseteq S\) is the admissible set, and \(H\subseteq A\)
is the selected target.  The associated reachability language is
\[
 \SEL(\mathfrak C)
 =
 \left\{w:\exists n\geq 0,\
 F^k(\iota(w))\in A\ (0\leq k\leq n),\
 F^n(\iota(w))\in H\right\}.
\]
Program A2 proves that this language is recursively enumerable under its
declared presentation and states the additional hypotheses needed for
decidability and undecidability.  Those results are imported here.

\begin{definition}[Claim ladder]
For one declared presentation, distinguish:
\begin{enumerate}[label=(C\arabic*)]
\item \emph{step computability}: \(F(s)\) can be computed when it is defined;
\item \emph{bounded reachability}: a target query is decided up to a supplied
finite horizon;
\item \emph{unbounded reachability}: membership in \(\SEL(\mathfrak C)\) is
decidable or undecidable;
\item \emph{complexity}: a total task has a lower bound in a specified machine
model and size measure;
\item \emph{computational irreducibility}: a separately defined class of
shortcut predictors fails on a specified family of instances.
\end{enumerate}
\end{definition}

The ladder is not a chain of implications.  In particular, a computable
one-step rule can have undecidable unbounded reachability, and an undecidable
language does not by itself state a runtime lower bound for a total decider,
because no such decider exists.

\begin{remark}[Robustness is another quantifier]
A robust event may mean that every initial point in a neighborhood reaches the
same target at the same step, that each point reaches it at some possibly
different step, or that a probability exceeds a threshold.  These are
different languages.  A paper must select one convention and provide an
effective representation of the quantified neighborhood.
\end{remark}

\section{Projection does not create undecidability}

The Projection--Admissibility Principle distinguishes descent from
recovery~\cite{NeroProjection}.  Information loss under descent is a structural
fact.  Computational universality is a property of a specified transition
system.  The following elementary example separates them.

\begin{proposition}[Decidable noninjective projection]\label{prop:counterexample}
There is a finite system with a noninjective projection, an absorbing stable
record basin, and selection reachability decidable in one step.
\end{proposition}

\begin{proof}
Let the source state space be \(Y=\{0,1\}^2\), the effective state space be
\(X=\{0,1\}\), and define
\[
 P(a,b)=a,\qquad G(a,b)=(0,b),\qquad F(a)=0.
\]
Then \(P\) is noninjective and
\[
  P\circ G=F\circ P.
\]
Take \(A=X\) and \(H=\{0\}\).  The target is absorbing.  Every effective state
is in \(H\) initially or enters \(H\) after one step, so the selection language
is the full input language and is decidable.  A finite product of copies with
componentwise dynamics gives the same conclusion for a local finite-lattice
model.
\end{proof}

The counterexample does not say that projection is irrelevant.  It says that
projection needs a fidelity contract before a computational property can pass
through it.

\subsection{The descent-transfer theorem}

Let
\[
\mathfrak C_Y=(S_Y,\iota_Y,G,A_Y,H_Y),
\qquad
\mathfrak C_X=(S_X,\iota_X,F,A_X,H_X)
\]
be computable admissible presentations.  Think of \(Y\) as an upper or source
description and \(X\) as an effective description.

\begin{definition}[Effective reachability descent]\label{def:descent}
An effective reachability descent from \(\mathfrak C_Y\) to
\(\mathfrak C_X\) on an invariant encoded set \(E\subseteq S_Y\) consists of
computable maps \(P:E\to S_X\), \(u\) on source inputs, and \(g\) on target
inputs such that:
\begin{enumerate}[label=(D\arabic*)]
\item \(\iota_Y(u(w))\in E\) and
\(P(\iota_Y(u(w)))=\iota_X(g(w))\);
\item \(G(E)\subseteq E\) wherever \(G\) is defined;
\item \(P(Gy)=F(Py)\) for every relevant \(y\in E\);
\item \(y\in A_Y\) if and only if \(Py\in A_X\), on the encoded orbits;
\item \(y\in H_Y\) if and only if \(Py\in H_X\), on the encoded orbits.
\end{enumerate}
\end{definition}

\begin{theorem}[Reachability transfer through projection]\label{thm:transfer}
If Definition~\ref{def:descent} holds, then for every declared input \(w\),
\[
 u(w)\in\SEL(\mathfrak C_Y)
 \quad\Longleftrightarrow\quad
 g(w)\in\SEL(\mathfrak C_X).
\]
Consequently, if a language \(L\) many-one reduces to
\(\SEL(\mathfrak C_Y)\) through \(u\), then
\[
 L\leq_m\SEL(\mathfrak C_X)
\]
through \(g\).
\end{theorem}

\begin{proof}
Condition (D1) identifies the initial states.  Conditions (D2) and (D3), by
induction on \(n\), give
\[
 P\!\left(G^n(\iota_Y(u(w)))\right)
 =
 F^n(\iota_X(g(w)))
\]
for every finite encoded prefix.  Condition (D4) preserves exactly the prefixes
that count as admissible, and (D5) preserves exactly the target hits.  The two
reachability statements are therefore equivalent.  Composition with the
source reduction proves the final claim.
\end{proof}

\begin{remark}
Noninjectivity of \(P\) away from \(E\) is compatible with the theorem.  What
matters is not invertibility of the whole projection but preservation of the
encoded orbit, the admissibility predicate, and the target predicate.  If
\(P\) collapses the halting and nonhalting codes together, the reduction is
destroyed rather than created.
\end{remark}

\section{Finite horizons and robust certificates}

Version 1 defined an instance-dependent finite admissible horizon and then
asserted that finiteness did not restore decidability.  That conclusion does
not follow.  A known computable horizon and decidable step predicates give a
finite search.  Conversely, merely naming a semantic \(N_{\max}\) does not
supply an algorithm for finding it.

For continuous or uncertain initial data, one also needs certified set
propagation.  The next result states a finite checker contract without
pretending that every nonlinear model automatically emits it.

\begin{definition}[Positive-margin enclosure packet]\label{def:packet}
Let \(X\) be a metric state space, \(K_0\subseteq A\) a compact preparation
set, \(F\) the selected step map, \(H\subseteq A\) an absorbing target, and
\(N<\infty\).  A positive-margin enclosure packet consists of sets
\(E_0,\ldots,E_N\) and machine-checkable rational bounds satisfying:
\begin{enumerate}[label=(V\arabic*)]
\item \(K_0\subseteq E_0\) and \(F(E_k)\subseteq E_{k+1}\) for \(k<N\);
\item \(\dist(E_k,X\setminus A)\geq a_k>0\) for every \(k\leq N\);
\item for every \(k\leq N\), exactly one certified row is supplied:
\[
\textsc{hit}: \dist(E_k,X\setminus H)\geq h_k>0,
\quad\text{or}\quad
\textsc{miss}: \dist(E_k,H)\geq m_k>0.
\]
\end{enumerate}
\end{definition}

\begin{proposition}[Finite-horizon robust decision]\label{prop:finite}
A valid packet from Definition~\ref{def:packet} decides whether all
preparations in \(K_0\) enter the absorbing target \(H\) by step \(N\).
\end{proposition}

\begin{proof}
Condition (V1) encloses every propagated preparation and (V2) keeps every
enclosed prefix admissible.  A \textsc{hit} row at step \(k\) places all
propagated states inside \(H\), so the answer is yes.  If every row through
\(N\) is \textsc{miss}, no propagated state has entered \(H\), so the answer is
no.  Absorption makes a hit by an earlier, preparation-dependent time visible
as a common hit at the final step.
\end{proof}

This proposition is intentionally a certificate theorem.  It does not say
that the required enclosures or strict distances are easy to compute.  If an
image touches \(\partial A\) or \(\partial H\), the positive-margin promise
fails and the packet is incomplete.  The correct output is then
\emph{unresolved at this precision}, not ``undecidable.''  Interval arithmetic,
semialgebraic decision procedures, abstract interpretation, or model-specific
barrier certificates may provide the rows in particular systems.

\section{The imported conditional A2 theorem}

The corrected Program A2 theorem uses one fixed universal two-counter machine
with configurations \(Q\times\N^2\)~\cite{Minsky1967,NeroA2}.  Its robust
uniform embedding contract has six indispensable rows:
\begin{center}
\begin{tabularx}{0.94\textwidth}{@{}>{\raggedright\arraybackslash}p{0.20\textwidth}Y@{}}
\toprule
Row & Required content \\
\midrule
Input code & A computable injection of machine configurations and a computable
map from finite machine inputs to physical initial data. \\
Step fidelity & One selected physical step implements increment,
decrement/zero-test, branch, and halt on every encoded run. \\
Unbounded persistence & Every finite prefix of every encoded run stays in the
domain and remains admissible; there is no fixed global counter cap. \\
Target fidelity & The selected event occurs exactly for halting configurations. \\
Uniformity & Dynamics, target, decoder, and protocol are fixed independently of
the input. \\
Robustness & Computable positive neighborhoods preserve labels, steps, and
target membership. \\
\bottomrule
\end{tabularx}
\end{center}

Program A2 proves the following imported implication:
\[
\boxed{
\text{complete robust uniform embedding}
\ \Longrightarrow\
\operatorname{HALT}_{\mathcal M}\leq_m\SEL(\mathfrak C)
}
\]
and therefore conditional undecidability of selection reachability.  The
proof is the standard finite-prefix simulation followed by target fidelity.
It is not repeated here, preserving theorem ownership in Program A2.

The force of the theorem lies in its antecedent.  Projection, non-Markovianity,
finite margins, locality, or the existence of records supplies none of the six
rows automatically.  Smooth and even analytic dynamical systems can be
constructed to simulate Turing machines~\cite{Branicky1995,GracaEtAl2005,
GracaEtAl2008,GracaZhong2023}; those results demonstrate possibility, not
genericity and not selection by MTT.

\section{A compact uniform-robustness obstruction}

The old record-register construction assumed a countable collection of stable
sites while also speaking as though one fixed finite capacity and one uniform
stability margin were sufficient.  Compactness exposes the conflict.

\begin{definition}[Decoded robust code]
Let \(K\) be a metric carrier and \(\mathcal C\) a set of symbolic
configurations.  A decoded robust code consists of points \(x_c\in K\),
radii \(\delta_c>0\), and a single-valued decoder
\[
 D:\bigcup_{c\in\mathcal C}B(x_c,\delta_c)\longrightarrow\mathcal C
\]
such that \(D(x)=c\) throughout \(B(x_c,\delta_c)\).
\end{definition}

\begin{theorem}[Compact uniform-margin obstruction]\label{thm:packing}
If \(K\) is compact, then for every \(r>0\) only finitely many distinct
configurations can satisfy \(\delta_c\geq r\).  In particular, an infinite
configuration set cannot have
\[
 \inf_{c\in\mathcal C}\delta_c>0.
\]
\end{theorem}

\begin{proof}
For distinct \(c,d\), the decoding neighborhoods are disjoint; otherwise the
single-valued decoder would assign two labels to one point.  Fix \(r>0\) and
consider all \(c\) with \(\delta_c\geq r\).  Their centers are \(r\)-separated:
if \(d(x_c,x_d)<r\), then \(x_d\) belongs to both
\(B(x_c,\delta_c)\) and \(B(x_d,\delta_d)\), a contradiction.

Compact metric spaces are totally bounded.  Cover \(K\) by finitely many balls
of radius \(r/2\).  Each such ball contains at most one center from an
\(r\)-separated set.  Hence only finitely many of the configurations can have
radius at least \(r\).  A uniform positive lower bound would contradict this
for an infinite configuration set.
\end{proof}

\begin{corollary}[Options for unbounded counters]\label{cor:options}
A robust encoding of two unbounded counters must use at least one of the
following:
\begin{enumerate}[label=(\roman*)]
\item decoding radii that approach zero along some configurations;
\item a carrier that is not totally bounded;
\item a uniformly specified family of growing carriers; or
\item a different encoding whose robustness is proved without one uniform
locally constant decoder margin.
\end{enumerate}
\end{corollary}

The theorem is deliberately limited.  It does not say that compact smooth
systems cannot simulate Turing machines.  Indeed, Gra\c{c}a and Zhong construct
a smooth simulation on a compact sphere~\cite{GracaZhong2023}.  It says that
one must inspect the exact encoding and robustness quantifiers: compactness is
incompatible with infinitely many disjoint code cells of one fixed positive
radius.  Shrinking cells or another coding mechanism evade that particular
obstruction but introduce a precision obligation that a physical claim must
address.

\section{Computable trajectories can have undecidable reachability}

The phrase ``no effective law exists'' in Version 1 was stronger than the
reduction supports.

\begin{proposition}[Finite prefixes versus unbounded reachability]
\label{prop:prefix}
There is a fixed total computable step map for which every finite orbit prefix
is computable, while target reachability over all time is undecidable.
\end{proposition}

\begin{proof}
Fix a universal Turing machine and encode its configurations as finite strings.
Let \(\tau\) be its computable one-step transition, extended by a self-loop on
halting configurations.  Given an input \(w\) and \(n\), the state
\(\tau^n(c_0(w))\) is computable by \(n\) iterations.  Thus every finite prefix
is computable.  If an algorithm decided whether some iterate entered the
halting set, it would decide the halting problem, contradicting
Turing's theorem~\cite{Turing1937}.
\end{proof}

The proposition locates the obstruction in the unbounded existential
quantifier over \(n\), not in the local update law.  Moore's dynamical-systems
constructions make the same distinction concrete in continuous
settings~\cite{Moore1990,Moore1991}.

\subsection{What ``computational irreducibility'' would still require}

Undecidable reachability is already a precise and strong statement.  It should
not be inflated into three other statements:
\begin{itemize}
\item It does not prove that every trajectory is difficult; many inputs may
halt immediately or enter a simple cycle.
\item It does not prove that each finite state requires step-by-step
simulation; a special family may admit a closed form or shortcut.
\item It does not provide a complexity lower bound without a total prediction
task, input-size convention, machine model, and comparison class.
\end{itemize}

A future MTT irreducibility theorem would therefore need to name a total
finite task, a family of selected instances, a resource measure, and a class of
allowed shortcut algorithms.  The conditional Program A2 theorem establishes
none of those complexity claims, nor does it purport to.

\section{What the external literature establishes}

The relevant literature supports a conditional, model-specific conclusion.
Hybrid and continuous systems can carry universal computation, but reachability
also has broad decidable subclasses.

\begin{center}
\begin{tabularx}{0.96\textwidth}{@{}p{0.23\textwidth}YY@{}}
\toprule
Result & What it establishes & What it does not establish \\
\midrule
Turing and Minsky
& Universal discrete machines and undecidable halting/reachability
& A physical realization or robustness under perturbations
\\
Moore
& Explicit dynamical systems whose long-run questions encode computation
& Generic undecidability of smooth dynamics
\\
Branicky
& Universal computation in specified hybrid and ODE constructions
& Universality from projection alone
\\
Henzinger et al.
& Precise decidable and undecidable reachability boundaries for hybrid automata
& One verdict for every hybrid system
\\
Gra\c{c}a et al.
& Explicit analytic or polynomial robust simulations, including bounded noise
& Selection of those dynamics by MTT
\\
Gra\c{c}a--Zhong
& Low-dimensional robust analytic simulation and a smooth compact-sphere
construction
& A uniform positive decoding radius for infinitely many labels
\\
\bottomrule
\end{tabularx}
\end{center}

The literature therefore strengthens the corrected paper in two ways.  It
shows that a robust embedding is mathematically plausible, and it shows why
the embedding must be constructed rather than inferred.  The decidability
frontier changes when guards, resets, dimensions, noise models, and encoding
conventions change~\cite{HenzingerEtAl1995}.

\section{Current MTT status}

MTT currently provides several relevant but separate layers:
\begin{enumerate}[label=(\arabic*)]
\item Foundations defines closure, admissibility, and typed claim
boundaries~\cite{NeroFoundation}.
\item Projection--Admissibility distinguishes descent, information loss, and
recovery~\cite{NeroProjection}.
\item Program A2 supplies the computable-presentation language, finite-depth
results, finite-capacity no-go, and the exact conditional universal-machine
reduction~\cite{NeroA2}.
\item The measurement paper distinguishes physical disturbance, outcome
completion, and record stabilization~\cite{NeroMeasurement}.
\item The conditional-classification paper prevents projection, event,
ultraviolet, and matter claims from being inferred across incomplete witness
axes~\cite{NeroClassification}.
\end{enumerate}

What is not currently selected is one physical realization satisfying all six
A2 embedding rows on the same source.  Finite basin models and finite event
logs can demonstrate bounded simulation.  They do not prove unbounded
admissible persistence.  A countable record register written into prose is not
a selected carrier.  A threshold metaphor is not an exact zero-test operator.
A stability margin for each finite experiment is not a uniform or controlled
family of margins over arbitrary counter values.

Accordingly, this paper makes no promotion of the physical MTT proof tier.  Its
new exact results are mathematical audit tools: the descent theorem identifies
what projection must preserve, and the compactness theorem identifies one
capacity/robustness obstruction.

\section{The certificate required for a physical theorem}

A future selected MTT undecidability claim should be accompanied by one
hash-addressed certificate with at least the following fields:
\begin{verbatim}
carrier_id, carrier_domain, state_representation,
input_encoder, configuration_decoder,
selected_step_operator, admissible_domain,
increment_map, decrement_map, zero_test,
halting_target, projection_map,
semiconjugacy_verifier, target_fidelity_verifier,
prefix_persistence_theorem, robustness_radii,
precision_model, source_hashes, proof_status
\end{verbatim}

The packet must answer concrete questions.
\begin{enumerate}[label=(R\arabic*)]
\item Is the carrier fixed for all inputs, or is there a uniform growing
family?
\item Can every finite machine prefix be represented without leaving the
selected admissible domain?
\item Are increment, decrement, and zero-test operations generated by one
fixed selected dynamics rather than an input-dependent controller?
\item Does the physical projection satisfy Theorem~\ref{thm:transfer} on the
encoded subset?
\item Are robustness radii uniform, shrinking, or controlled by a proved
precision law?
\item Does the target event represent halting in both directions?
\end{enumerate}

Any failed row is informative.  Failure of unbounded persistence leaves a
finite simulator.  Failure of target fidelity invalidates the reduction.
Failure of robustness leaves an exact-code construction that may be physically
unstable.  Failure of semiconjugacy means the projected event language need not
inherit the source computation.

\section{Completion and falsifiability}

The corrected framework can fail cleanly.
\begin{enumerate}[label=(F\arabic*)]
\item A proposed projection identifies a halting and nonhalting encoded orbit,
violating target fidelity.
\item The counter code reaches a fixed capacity or an admissibility boundary.
\item Zero testing requires precision smaller than the declared physical
tolerance.
\item The protocol or transition rule depends on the encoded program, so the
claimed fixed universal dynamics is not uniform.
\item Robust neighborhoods overlap while carrying distinct labels.
\item A finite-horizon claim lacks a positive-margin enclosure packet.
\item A claimed complexity or irreducibility theorem lacks a total task and
resource model.
\end{enumerate}

Conversely, completing the six A2 rows and the descent contract would justify a
substantial theorem: one selected MTT event-reachability language would be
undecidable.  It still would not imply that every MTT process is universal or
that all finite predictions require literal simulation.

\section{Conclusion}

Projection is not an engine of undecidability.  It can erase a computation,
preserve one, or transfer its reachability language, depending on the maps and
predicates it respects.  Stable records provide possible symbolic carriers,
but they do not supply unbounded counters, zero tests, or universal dynamics.
Finite admissibility can support verified finite simulations; it cannot be
silently promoted to unlimited machine execution.

The corrected result is therefore both smaller and stronger.  Program A2 owns
the exact conditional reduction.  This paper proves the descent contract
needed to apply it through a projection, gives a decidable projection
counterexample, supplies a finite-horizon certificate theorem, and proves a
compact uniform-margin obstruction.  It also makes explicit that computable
finite trajectories and undecidable unbounded reachability coexist.

For MTT, the frontier is now unambiguous: construct one selected physical
carrier and operator satisfying the complete robust embedding and descent
certificate, or retain the valuable but finite simulation tier.  Until that
construction exists, selection undecidability is a conditional mathematical
possibility, not a derived universal property of reality.

\bibliographystyle{plain}
\bibliography{main}

\end{document}
